\section{Introduction}
\label{sec:intro}

Brazil's general procurement statute (Lei 8.666/1993, Art.\ 22,
\S\,3) requires at least three participants in sealed-bid
(\emph{convite}) tenders. The rule exists to prevent collusion: more
bidders should mean more competition and lower prices. This paper
shows that the rule can achieve the opposite. When genuine
competitors are scarce, bid-rigging cartels deploy \emph{cover
bidders}---firms that submit deliberately non-competitive bids solely
to satisfy the participation floor---turning a pro-competitive
safeguard into a subsidy for the organizational costs of collusion.

The mechanism is straightforward. A cartel that controls the winning
bid needs the tender to clear the three-bidder threshold; otherwise
the procuring unit must cancel and re-advertise, attracting
unwanted competition. So the cartel assigns affiliated firms to
submit bids just high enough to lose. These firms---which we call
\emph{frequent losers} (FL)---accumulate dozens of participations
with a zero win rate, a pattern that is difficult to rationalize as
competitive entry when bidding is costly
\citep{marshall2012economics, asker2010leniency}. The minimum-bidder
rule does not create cartels, but it creates a task that cartels
must perform, and that task leaves a detectable trace.

We study this mechanism using 4.5 million tender-items and 40 million
bids from S\~ao Paulo's Bolsa Eletr\^onica de Compras (BEC,
2009--2019). The BEC platform offers a natural institutional
comparison: \emph{convite} tenders are subject to the three-bidder
rule, while \emph{preg\~ao} (electronic reverse auction) tenders,
introduced by Lei 10.520/2002, are not. If the minimum-bidder rule
facilitates cover bidding, its effects should be most visible in
convite---but the cover-bidding infrastructure, once built, can be
redeployed in preg\~ao as well, where it inflates prices without any
institutional mandate.

Three empirical findings support this account. \emph{First}, among
FL-flagged convite tenders, 39.2\% have fewer than three genuine
competitors, forcing the cartel to deploy cover firms to clear the
threshold. In 23.5\% of cases, every participant is a cover
firm---the rule produced zero genuine competition at positive
administrative cost. These figures describe a corner solution: the
rule binds, and the cartel fills the gap.

\emph{Second}, the price effect of cover-bidder presence is
asymmetric across regimes. FL-flagged convite tenders show 3.8\%
higher prices in within-item comparisons; FL-flagged preg\~ao tenders
show 9.3\%. The pattern is consistent with the minimum-bidder
constraint \emph{diluting} rather than generating the collusive
signal: convite forces some cover bidding in markets where
competition is already thin, compressing the coefficient toward zero,
while preg\~ao captures only voluntary deployment where the cartel
expects it to be profitable. The 7.6 percentage-point convite premium
in the subsample where the rule does not bind (tenders with $\geq 3$
genuine bidders) reinforces this reading.

\emph{Third}, the price effect varies 12.5-fold across purchasing
units (PBUs), concentrating where oversight capacity is weakest.
Small PBUs (bottom quartile by annual procurement volume) show the
largest FL price gaps; large PBUs show economically negligible
effects. This is consistent with models where detection probability
constrains cartel behavior \citep{harrington2008detecting}: where
oversight is credible, the rule's collusion-facilitating potential is
disciplined.

Back-of-the-envelope welfare calculations suggest the minimum-bidder
rule costs R\$74--211 million over the sample period---a costless
institutional reform that requires only legislative amendment, not
new enforcement capacity. Replacing convite with preg\~ao as the
default low-value modality, as several Brazilian states have done
since 2019, would eliminate the rule's bite without reducing
procedural safeguards.

\subsection{Contributions and Related Literature}

This paper contributes to three literatures.

\emph{First}, it adds to the law-and-economics literature on how
procurement rules shape competitive outcomes. Minimum-bidder
requirements are common across jurisdictions---the European Union's
Directive 2014/24/EU, the World Bank's procurement guidelines, and
similar statutes in India, South Africa, and several U.S.\
states all impose participation floors in some procurement
modalities. The unintended consequences of such rules have been
discussed theoretically \citep{decarolis2014awarding}, but
direct empirical evidence on how they facilitate collusion is
scarce. We provide it.

\emph{Second}, it contributes to the empirical literature on
bid-rigging detection by documenting the institutional channel
through which cover bidders operate. Existing screens focus on
bid-level patterns
\citep{bajari2003deciding, imhof2019detecting, chassang2022robust};
we show that a simple participation anomaly---zero wins, abnormally
many tenders---captures information orthogonal to these tools and
is especially informative where the institutional design creates
demand for cover bidding.%
\footnote{The frequent-loser screen and its validation are
developed in a companion paper \citep[Paper~1]{genicolomartins2026screening}.
This paper takes the screen as given and studies its institutional
determinants.}

\emph{Third}, it speaks to the policy debate on procurement reform
in Brazil and comparable jurisdictions. Lei 14.133/2021 (the
``New Procurement Law'') replaced Lei 8.666 but retained
minimum-bidder provisions for certain modalities. Our findings
suggest that these provisions warrant reconsideration: the
participation floor is binding precisely where it does the most
harm.

The rest of the paper proceeds as follows. Section~\ref{sec:institutional}
describes the institutional background: the convite and preg\~ao
modalities, the minimum-bidder rule, and the BEC platform.
Section~\ref{sec:framework} develops the theoretical framework
linking the rule to cover-bidder deployment.
Section~\ref{sec:data} describes the data and the frequent-loser
classification. Section~\ref{sec:empirical} presents the empirical
strategy. Section~\ref{sec:results} reports the results.
Section~\ref{sec:policy} discusses policy implications and reform
options. Section~\ref{sec:conclusion} concludes.
